\exercisegroup

\begin{enumerate}
\def\labelenumi{\arabic{enumi})}
\item
  设 \(\mathrm{F}, \mathrm{G}\) 是处于对偶中的两个向量空间，且 \(\sigma(\mathrm{F}, \mathrm{G})\) 是 Hausdorff 的。证明：如果 \(\mathcal{T}\) 是一个与 \(\mathrm{F}\) 的向量空间结构相容的 Hausdorff 拓扑，并且粗于 \(\sigma(\mathrm{F}, \mathrm{G})\) （但先验地未必局部凸），那么 \(\mathcal{T} = \sigma(\mathrm{F}, \mathrm{G}_1)\) ，其中 \(\mathrm{G}_1\) 是 G 的一个向量子空间，它在拓扑 \(\sigma(\mathrm{G},\mathrm{F})\) 中稠密。（在 F 上考虑这样的局部凸拓扑 \(T_{1}\) ：其 0 的一个基本邻域系由 T 中那些闭的、凸的、均衡的且为 \(\sigma(\mathrm{F},\mathrm{G})\) 之 0 邻域的集组成。）由此推出：如果 \(T_{0}\) 是向量空间 E 上的一个 Hausdorff 局部凸拓扑，它在 E 上的 Hausdorff 局部凸拓扑所成之集中是极小的（II, 第 85 页, 习题 13），那么它在由与 E 的向量空间结构相容的 Hausdorff 拓扑（无论局部凸与否）所成之集中也是极小的。
\item
  在 \(\mathbf{R}^n\) 中，设 \((\mathrm{C}_i)_{1\leqslant i\leqslant m}\) 是由 \(m\geqslant n + 1\) 个以 0 为顶点的凸锥组成的族；证明：如果对这些锥中任何 \(n + 1\) 个锥，都存在一个过 0 的超平面 \(\mathbf{H}\) ，使这些锥位于 \(\mathbf{H}\) 的同侧，那么存在一个超平面 \(\mathrm{H}_0\) ，使所有的锥 \(\mathrm{C}_i(1\leqslant i\leqslant m)\) 都位于 \(\mathrm{H}_0\) 的同侧（参见 II, 第 68 页, 习题 21, a)）。
\item
  在 \(\mathbf{R}^n\) 中，设 \((\mathrm{D}_i)_{1\leqslant i\leqslant m}\) 是由 \(m\geqslant 2n\) 个由过 0 的超平面所确定的闭半空间组成的族。证明：如果对这些半空间中任何 \(2n\) 个，其交中存在一个 \(\neq 0\) 的点，那么存在一个 \(\neq 0\) 的点，它属于所有 \(\mathrm{D}_i(1\leqslant i\leqslant m)\) 的交（参见 II, 第 66 页, 习题 10）。
\item
  设 S、T 是 \(R^{n}\) 中两个没有公共点的有限集，且它们的并至少含 \(2n + 2\) 个点。那么，存在分离 S 与 T 的超平面，其充要条件是：对每个有限集 \(F \subset S \cup T\) （由 \(2n + 2\) 个点组成），都存在分离 \(F \cap S\) 与 \(F \cap T\) 的超平面（利用习题 3 以及 II, 第 78 页, 习题 15 的方法）。
\item
  设 E 是 \(R^{n}\) 上二次型所成的向量空间，它等同于 R 上 n 阶方阵所成空间 \(\mathbf{M}_{n}(\mathbf{R})\) 中由对称方阵组成的向量子空间。我们给 \(\mathbf{M}_{n}(\mathbf{R})\) 装备内积 \(\operatorname{Tr}(^{t}\mathbf{X},\mathbf{Y})\) ，借此可把它与其对偶等同起来；对 E 亦同样处理。
\end{enumerate}

\begin{enumerate}
\def\labelenumi{\alph{enumi})}
\item
  设 \(\mathbf{P} \subset \mathbf{E}\) 是其矩阵的元素全 \(\geqslant 0\) 的二次型所成之集，\(\mathbf{S} \subset \mathbf{E}\) 是 \(\mathbf{R}^n\) 中正二次型所成之集。证明：\(\mathbf{P} = \mathbf{P}^\circ\) ，且 \(\mathbf{S} = \mathbf{S}^\circ\) 。
\item
  设 B 是 \(\mathbf{R}^n\) 上形如 \(\sum_{j=1}^{m} x_j'^2\) （对某个 \(m\) ）的二次型所成之集，其中 \(x_{j}^{\prime}\) 是取值 \(\geqslant 0\) 的线性型——对一切向量 \(x = (x_{i})_{1\leqslant i\leqslant n}\) （其坐标 \(x_{i}\) 全 \(\geqslant 0\) ）而言；设 \(C\) 是对所有向量取值 \(\geqslant 0\) 的二次型所成之集，这里向量 \(x = (x_{i})\) 的坐标 \(x_{i}\) 全 \(\geqslant 0\) 。证明：\(B = C^{\circ}\) 且 \(C = B^{\circ}\) （证明 \(B\) 是闭的：证明 \(B\) 的每个元素都可写成形式 \(\sum_{j=1}^{m} x_{j}^{\prime 2}\) ，其中 \(x_{j}^{\prime}\) 对一切 \(x\) （坐标 \(\geqslant 0\) ）为正，且 \(m \leqslant 2^{n}\) ）。
\end{enumerate}

\begin{enumerate}
\def\labelenumi{\arabic{enumi})}
\setcounter{enumi}{5}
\tightlist
\item
  设 F、G 是处于分离对偶中的两个向量空间，A 是 F 中一个弱紧凸集。设 C 是 G 中一个以 0 为顶点且弱闭的凸锥。假设对一切 \(y \in \mathbf{C}\) ，存在 \(x \in \mathbf{A}\) 使 \(\langle x, y \rangle \geqslant 0\) 。证明：存在 \(x_0 \in \mathbf{A}\) ，使 \(\langle x_0, y \rangle \geqslant 0\) 对一切 \(y \in \mathbf{C}\) 成立（对 A 与 \(\mathbf{C}^\circ\) 应用 II, 第 38 页, prop. 4）。
\end{enumerate}

¶ 7) a) 设 F、G 是处于分离对偶中的两个向量空间，C 是 F 中一个弱闭凸锥。设 M 是 G 的一个有限维向量子空间。证明：或者存在 \(y_0 \in \mathbf{C}\) ，使 \(y_0 \in \mathbf{M}^\circ\) 且 \(y_0 \neq 0\) ，或者存在 \(z_0 \in \mathbf{M}\) ，使 \(z_0 \in \mathbf{C}^\circ\) 且 \(z_0 \neq 0\) （对 M 的维数作归纳论证）。如果 C 不含任何直线，而上述两个性质同时成立，证明 \(z_0\) 不可能是 \(\mathbf{C}^\circ\) 的代数内点。b) 设 \((a_{ij}), (b_{ij})\) 是两个具有 \(n\) 行 \(m\) 列实元素的矩阵，且 \(a_{ij} > 0\) 对每一对 \((i,j)\) 成立。证明：存在唯一的值 \(\lambda \in \mathbf{R}\) ，使得存在两个向量 \(x = (x_j) \in \mathbf{R}^m\) 、\(y = (y_i) \in \mathbf{R}^n\) ，满足关系 \(x \neq 0, y \neq 0, x_j \geqslant 0, y_i \geqslant 0\) （对一切 \(i,j\) ），并且最后满足

\[
\lambda \sum_ {j = 1} ^ {m} a _ {i j} x _ {j} \geqslant \sum_ {j = 1} ^ {m} b _ {i j} x _ {j} \quad \text { for } \quad 1 \leqslant i \leqslant n\tag{1}
\]

\[
\lambda \sum_ {i = 1} ^ {n} a _ {i j} y _ {i} \leqslant \sum_ {i = 1} ^ {n} b _ {i j} y _ {i} \quad \text { for } \quad 1 \leqslant j \leqslant m\tag{2}.
\]

（令 \(c_{ij} = \lambda a_{ij} - b_{ij}\) （对 \(1 \leqslant i \leqslant n\) ，\(1 \leqslant j \leqslant m\) ），并令 \(c_{n + i,j} = \delta_{ij}\) （Kronecker 指标，对 \(1 \leqslant i \leqslant m\) ），证明问题归结为求一个非零向量 \(x \in \mathbf{R}^m\) 与一个非零向量 \(z = (z_i) \in \mathbf{R}^{n + m}\) ，使它们满足关系

\[
\sum_ {j = 1} ^ {m} c _ {i j} x _ {j} \geqslant 0 \quad \text { for } \quad 1 \leqslant i \leqslant n + m\tag{3}
\]

\[
\sum_ {i = 1} ^ {n + m} c _ {i j} z _ {i} = 0 \quad \text { for } \quad 1 \leqslant j \leqslant m\tag{4}
\]

以及 \(z_{i} \geqslant 0\) （对 \(1 \leqslant i \leqslant n + m\) ）。注意：如果 (3) 在某个值 \(\lambda_0\) （\(\lambda\) 的一个值）处有解，那么它对 \(\lambda \geqslant \lambda_0\) 也有解；如果 (4) 对 \(\lambda_0\) 有解，那么它对 \(\lambda \leqslant \lambda_0\) 也有解。最后利用 \(a)\) 。

\begin{enumerate}
\def\labelenumi{\arabic{enumi})}
\setcounter{enumi}{7}
\tightlist
\item
  设 \(\mathbf{T}\) 是紧空间，\(\mathbf{L}\) 是 \(\mathcal{C}(\mathbf{T};\mathbf{R})\) 的一个维数为 \(r\) 的向量子空间；给 \(\mathbf{L}\) 装备由 \(\mathcal{C}(\mathbf{T};\mathbf{R})\) 的范数诱导的范数，并给其对偶 \(\mathbf{L}^*\) 装备范数 \(\| x'\| = \sup_{\| x\| \leqslant 1} \langle x, x' \rangle\) ，于是，若 \(\mathbf{B}\) 是球 \(\| x\| \leqslant 1\) （在 \(\mathbf{L}\) 中），则 \(\mathbf{B}^\circ\) 是球 \(\| x'\| \leqslant 1\) （在 \(\mathbf{L}^*\) 中）。
\end{enumerate}

\begin{enumerate}
\def\labelenumi{\alph{enumi})}
\tightlist
\item
  对每个 \(t \in \mathbf{T}\) ，用 \(e_t'\) 表示线性型 \(x \mapsto x(t)\) （\(\mathbf{L}\) 上的）。证明：\(\mathbf{B}^\circ\) 是诸 \(\pm e_t'\) 所成之集的凸包，其中 \(t\) 取遍 \(\mathbf{T}\) （利用 II, 第 38 页, prop. 4 得出矛盾）。
\end{enumerate}

由此推出：每个线性型 \(x' \in \mathbf{L}^*\) （满足 \(\| x' \| = 1\) ）都可写成 \(x' = \sum_{i=1}^{r} \lambda_i e_{t_i}'\) ，其中

诸 \(t_i\) 是 \(r\) 个属于 \(\mathbf{T}\) 的点，诸 \(\lambda_i\) 是实数，满足 \(\sum_{i=1}^{r} |\lambda_i| = 1\) （参见 II, 第 66 页, 习题 9, a)）。

\begin{enumerate}
\def\labelenumi{\alph{enumi})}
\setcounter{enumi}{1}
\tightlist
\item
  对每个 \(y \in \mathcal{C}(\mathbf{T}; \mathbf{R})\) ，存在唯一的 \(x \in \mathbf{L}\) 使 \(\| y - x \| = d(y, \mathbf{L})\) ，其充要条件是：对每个非零的 \(z \in \mathbf{L}\) ，至多存在 \(r - 1\) 个互异的点 \(t_i \in \mathbf{T}\) 使 \(z(t_i) = 0\) （Haar 定理）。（为证此条件是充分的，首先注意它等价于说：对 \(r\) 个互异的点 \(t_i \in \mathbf{T}\) （ \(1 \leqslant i \leqslant r\) ），诸 \(e_{t_i}'\) 在 \(\mathbf{L}^*\) 中线性无关。然后通过假设结论不真并得出矛盾来论证。如果存在两个互异的点 \(x', x''\) （属于 \(\mathbf{L}\) ）使 \(\| y - x' \| = \| y - x'' \| = d(y, \mathbf{L})\) ，那么存在 \(x_0 \in \mathbf{L}\) 与 \(z \in \mathbf{L}\) ，使得对一切充分小的实数 \(\lambda\) ，有 \(\| y - (x_0 + \lambda z) \| = d(y, \mathbf{L})\) 。把 \(a\) ) 的最后一部分应用于子空间 \(\mathbf{L} \oplus \mathbf{R}y\) （\(\mathcal{C}(\mathbf{T}; \mathbf{R})\) 的）以及该空间上一个在 \(\mathbf{L}\) 上为零的适当线性型。为看出该条件是必要的，注意若它不真，则存在 \(r\) 个互异的点 \(t_i \in \mathbf{T}\) （ \(1 \leqslant i \leqslant r\) ），使诸 \(e_{t_i}'\) 线性相关，且存在一个非零的函数 \(z \in \mathbf{L}\) ，它在诸点 \(t_i\) 处为零。若 \(\alpha_i (1 \leqslant i \leqslant r)\) 是不全为零的数，使 \(\sum_{i=1}^{r} \alpha_i e_{t_i}' = 0\) ，则考虑一个函数 \(w \in \mathcal{C}(\mathbf{T}; \mathbf{R})\) ，使 \(\| w \| = 1\) ，\(w(t_i) = \text{sgn}(\alpha_i)\) （对 \(1 \leqslant i \leqslant r\) ），以及函数 \(y = w(1 - |\beta z|)\) ，其中 \(|\beta|\) 充分小且 \(\neq 0\) 。）
\item
  设 \(T\) 是 \(R\) 中的一个紧区间，且 \(L\) 满足 Haar 定理的条件；设 \((t_i)_{1 \leqslant i \leqslant r+1}\) 是由 \(r + 1\) 个属于 \(T\) 的点组成的严格递增序列；则存在 \(r + 1\) 个非零实数 \(\lambda_i\) ，使 \(\sum_{i=1}^{r+1} \lambda_i e_{t_i}' = 0\) 。证明此时 \(\text{sgn}(\lambda_i) \text{sgn}(\lambda_{i+1}) = -1\) 对 \(1 \leqslant i \leqslant r\) 成立。（分别考虑 \(r = 1\) 与 \(r > 1\) 两种情形。在第二种情形，相反地假设：对某个指标 \(i \leqslant r - 1\) ，数 \(\lambda_i\) 与 \(\lambda_{i-1}\) 或与 \(\lambda_{i+1}\) 同号，且 \(\lambda_{i-1}\) 与 \(\lambda_{i+1}\) 异号。若例如 \(\lambda_i > 0\) ，取 \(\alpha_{i-1} > 0\) 、\(\alpha_{i+1} > 0\) 使 \(\alpha_{i-1} \lambda_{i-1} + \alpha_{i+1} \lambda_{i+1} = 0\) ，再取 \(z \in L\) 使 \(z(t_{i-1}) = \alpha_{i-1}\) 、\(z(t_{i+1}) = \alpha_{i+1}\) ，且 \(z(t_j) = 0\) ，其中 \(j\) 异于 \(i - 1, i\) 与 \(i + 1\) 。推出 \(z(t_i) < 0\) ，并证明这与所给假设矛盾。）
\end{enumerate}

\(d)\) ）在与 \(c\) ) 相同的假设与记号下，设存在 \(y \in \mathcal{C}(\mathrm{T}; \mathbf{R})\) 与 \(z \in \mathrm{L}\) ，使 \(y(t_i) - z(t_i) = (-1)^i \alpha_i\) ，其中 \(\alpha_i > 0\) （对 \(1 \leqslant i \leqslant r + 1\) ）。证明此时有 \(d(y, \mathrm{L}) \geqslant \inf \alpha_i\) 。（把 \(a\) ) 应用于 \(\mathrm{L} \oplus \mathbf{R}y\) ，并用 \(c\) 。）

\begin{enumerate}
\def\labelenumi{\alph{enumi})}
\setcounter{enumi}{4}
\tightlist
\item
  在 \(c\) ) 的假设下，设 \(y \in \mathcal{C}(\mathbf{T};\mathbf{R})\) ，并设 \(z\) 是 \(\mathbf{L}\) 中使 \(\| y - z \| = d(y,\mathbf{L})\) 的唯一点。证明：存在一个严格递增序列 \((t_i)_{1\leqslant i\leqslant r + 1}\) （在 \(\mathbf{T}\) 中），使得
\end{enumerate}

\[
y (t _ {i}) - z (t _ {i}) = (- 1) ^ {i} \varepsilon \| y - z \|
\]

其中 \(\varepsilon = \pm 1\) 。反之，若 \(z\) 具有此性质，则 \(z\) 是 \(\mathbf{L}\) 中使 \(\| y - z \| = d(y, \mathbf{L})\) 的唯一点。（利用 \(c\) ) 与 \(d\) ）。考察 \(\mathbf{T}\) 是 \(\mathbf{R}\) 的一个区间且 \(\mathbf{L}\) 是限制到 \(\mathbf{T}\) 上的、次数 \(< r\) 的多项式所成之集的情形（\(Tchebycheffs\) 定理）。

\begin{enumerate}
\def\labelenumi{\arabic{enumi})}
\setcounter{enumi}{8}
\tightlist
\item
  设 \(\mathbf{F}\) 与 \(\mathbf{G}\) 是处于分离对偶中的两个向量空间，\(\mathbf{A}\) 是 \(\mathbf{F}\) 的一个包含 0 的凸子集。对每个 \(y \in \mathbf{G}\) ，令
\end{enumerate}

\[
\mathrm{H} _ {\mathbf {A}} (y) = \sup _ {x \in \mathbf {A}} (- \langle x, y \rangle),
\]

于是 \(0 \leqslant \mathrm{H}_{\mathbf{A}}(y) \leqslant +\infty\) ；我们称 \(\mathrm{H}_{\mathbf{A}}\) 为 \(\mathbf{A}\) 的支撑函数。

\begin{enumerate}
\def\labelenumi{\alph{enumi})}
\item
  证明：\(\mathbf{H}_{\mathrm{A}}\) 是 \(\mathbf{A}^{\circ}\) 的度规（II, 第 20 页）。
\item
  如果 A 是弱紧的，那么对一切 \(y \in \mathbf{G}\) ，方程为 \(\langle x, y \rangle = \mathrm{H}_{\mathbf{A}}(y)\) 的超平面是 A 的一个支撑超平面。
\item
  \(\mathbf{H}_{\mathrm{A}}\) 对拓扑 \(\sigma (\mathbf{G},\mathbf{F})\) 而言有限且连续，其充要条件是 A 有限维且有界（在它所生成的有限维向量子空间中）。
\end{enumerate}

\(d\) ) 设 \(\mathbf{A}_i(1\leqslant i\leqslant p)\) 是 F 中包含 0 的凸集，\(\lambda_{i}\) 是 \(\geqslant 0(1\leqslant i\leqslant p)\) 的实数；证明凸集 \(\mathbf{A} = \sum_{i}\lambda_{i}\mathbf{A}_{i}\) 的支撑函数是 \(\mathrm{H}_{\mathrm{A}} = \sum_{i}\lambda_{i}\mathrm{H}_{\mathrm{A}_{i}}\) 。如果 \(y\in \mathbf{G}\) 使得

交集 \(C_{i}\) （即 \(A_{i}\) 与超平面 \(\langle x, y \rangle = H_{A}(y)\) 之交）对 \(1 \leqslant i \leqslant p\) 非空，证明 A 与方程为 \(\langle x, y \rangle = H_{A}(y)\) 的超平面之交就是集合 \(\sum \lambda_{i}C_{i}\) 。

\begin{enumerate}
\def\labelenumi{\alph{enumi})}
\setcounter{enumi}{4}
\item
  设 \(\mathbf{A}\) 局部紧且不含任何直线。那么，由 \(\{0\}\) 与诸 \(y \neq 0\) （它们满足 \(\mathrm{H}_{\mathrm{A}}(y) = +\infty\) ，\(\mathrm{H}_{\mathrm{A}}(-y) \neq +\infty\) ）所成之集的并，是渐近锥 \(\mathrm{C}_{\mathrm{A}}\) 的极锥（考察 \(\mathbf{A}\) 是 \(\{0\}\) 与一条半直线之凸包的情形）。
\item
  设 F 有限维。证明：A 是抛物的（II, 第 67 页, 习题 17），其充要条件是 \(H_{A}\) 是 G 到 \(\overline{R}\) 中的一个连续映射（如果存在一条平行于 \(C_{A}\) 的某条半直线且不与 A 相交的直线，注意存在一个分离该半直线与 A 的超平面）。
\end{enumerate}

\begin{enumerate}
\def\labelenumi{\arabic{enumi})}
\setcounter{enumi}{9}
\tightlist
\item
  对 \(E = R^{n}\) 中每个包含 0 的紧凸集 A，我们借助 E 与 E* 之间的对偶，使它的支撑函数 \(H_{A}\) 与之对应：\(E^{*} : H_{A}\) 属于空间 \(\mathcal{C}(E^{*}; \mathbf{R})\) ，即 \(E^{*}\) 上连续实值函数所成之空间。我们给空间 \(\mathcal{C}(E^{*}; \mathbf{R})\) 装备紧收敛的一致结构，并给 E 中包含 0 的紧凸集所成之集 \(\Re_{0}^{\prime}(E)\) 装备 II, 第 71 页, 习题 39 中定义的一致结构。证明：\(A \mapsto H_{A}\) 是 \(\Re_{0}^{\prime}(E)\) 到 \(\mathcal{C}(E^{*}; \mathbf{R})\) 的一个一致子空间上的同构。
\end{enumerate}

由此推出：映射 \(\mathbf{A} \mapsto \mathbf{A}^{\circ}\) （由集 \(\Re_0(\mathbf{E})\) ——即 \(\mathbf{E}\) 中以 0 为内点的紧凸集全体——到集 \(\Re_0(\mathbf{E}^*)\) 上）对这两个空间的一致结构而言是同构（参见 II, 第 71 页, 习题 39）。

\begin{enumerate}
\def\labelenumi{\arabic{enumi})}
\setcounter{enumi}{10}
\tightlist
\item
  设 \(\mathrm{F}, \mathrm{G}\) 是处于分离对偶中的两个向量空间。\(\mathfrak{U}\) 是 \(\mathrm{F}\) 上的一个超滤子，它弱收敛于一点 \(x_0\) 的充要条件是：\(x_0\) 属于所有属于 \(\mathfrak{U}\) 的弱闭凸集之交（注意：若 \(x_0\) 是该交中一点而非 \(\mathfrak{U}\) 的聚点，则存在一个属于 \(\mathfrak{U}\) 且不含 \(x_0\) 的闭半空间）。
\end{enumerate}

由这一结果推出：为使 F 的点列 \((x_{n})\) 弱收敛于一点 a，必要且充分的是：a 属于由该序列的无穷多项所成之集的诸弱闭凸包（利用 GT, I, § 6.4 的 prop. 7）。

\begin{enumerate}
\def\labelenumi{\arabic{enumi})}
\setcounter{enumi}{11}
\tightlist
\item
  \begin{enumerate}
  \def\labelenumii{\alph{enumii})}
  \tightlist
  \item
    设 E 是一个向量空间，\((\mathrm{E}_{\alpha})_{\alpha\in\mathrm{A}}\) 是 E 的子空间的一个递增定向族，其并为 E；设每个 \(E_{\alpha}\) 上带有一个局部凸拓扑 \(T_{\alpha}\) ，使得对 \(\alpha\leqslant\beta\) ，典范单射 \(E_{\alpha}\to E_{\beta}\) 连续。设 T 是 E 上的拓扑，它是诸 \(T_{\alpha}\) 的归纳极限（II, 第 29 页, 例 II）；证明：E 的对偶 \(E'\) （对 T 而言）带拓扑 \(\sigma(\mathrm{E}',\mathrm{E})\) 可以典范地等同于诸对偶 \(E'_{\alpha}\) 带拓扑 \(\sigma(\mathrm{E}'_{\alpha},\mathrm{E}_{\alpha})\) 的投影极限。
  \end{enumerate}
\end{enumerate}

\begin{enumerate}
\def\labelenumi{\alph{enumi})}
\setcounter{enumi}{1}
\tightlist
\item
  设 \((\mathbf{X}_{\alpha},\phi_{\alpha \beta})\) 是与指标的一个定向集 A 相对应的非空集合的投影体系，使得诸 \(\phi_{\alpha \beta}\) 都是满射且 \(\varprojlim \mathbf{X}_{\alpha} = \emptyset\) \((\mathbf{S},\mathbf{III},\S 7,\text{exerc. }4)\) 。令 \(\mathbf{F}_{\alpha} = \mathbf{R}^{(\mathbf{X}_{\alpha})}\) ，并用 \(f_{\alpha \beta}:\mathrm{F}_{\beta}\to \mathrm{F}_{\alpha}\) （对 \(\alpha \leqslant \beta\) ）表示从 \(\phi_{\alpha \beta}\) 典范地导出的线性映射（A, II, § 1.11, cor. 1）。如果我们给每个 \(\mathbf{F}_{\alpha}\) 装备其因子的直和拓扑，那么对偶 \(\mathrm{E}_{\alpha} = \mathrm{F}_{\alpha}'\) （\(\mathrm{F}_{\alpha}\) 的对偶）带弱拓扑 \(\sigma(\mathrm{E}_{\alpha}, \mathrm{F}_{\alpha})\) 可以等同于乘积空间 \(\mathbf{R}^{\mathbf{X}_{\alpha}}\) ，且 \({}^{t}f_{\alpha\beta}\) 是 \(\mathrm{E}_{\alpha}\) 到 \(\mathrm{E}_{\beta}\) 的一个闭子空间上的同构，该子空间在 \(\mathrm{E}_{\beta}\) 中有拓扑补。证明：在 \(\mathrm{E} = \varinjlim \mathrm{E}_{\alpha}\) 上（对诸 \({}^{t}f_{\alpha\beta}\) 而言），诸 \(\mathrm{E}_{\alpha}\) 的拓扑的归纳极限是最粗的拓扑（从而非 Hausdorff）（利用 \(a\) 并注意到 \(\varprojlim \mathrm{F}_{\alpha} = \{0\}\) ）。
\end{enumerate}

\begin{enumerate}
\def\labelenumi{\arabic{enumi})}
\setcounter{enumi}{12}
\tightlist
\item
  设 E 是一个向量空间。我们称 E 上的一个 Hausdorff 局部凸拓扑 T 是极小的（并称带 T 的 E 是极小型空间），如果在 E 上不存在严格粗于 T 的 Hausdorff 局部凸拓扑（参见 II, 第 81 页, 习题 1）。
\end{enumerate}

\begin{enumerate}
\def\labelenumi{\alph{enumi})}
\item
  设 \(\mathcal{T}\) 是 E 上的一个极小拓扑，并设 \(E'\) 是 E 的对偶（当 E 带拓扑 \(\mathcal{T}\) 时）；证明 \(\mathcal{T} = \sigma(E, E')\) 且 \(E = E'^*\) （注意：利用 II, 第 43 页的推论 3，在 \(E'\) 中对拓扑 \(\sigma(E', E)\) 不可能存在处处稠密的超平面）。由此推出极小型空间是直线的乘积。
\item
  证明：在 Hausdorff 局部凸空间 F 中，每个极小型的子空间 E 都有拓扑补，从而特别地是闭的（利用 a) 以及 Hahn--Banach 定理，把 E 到自身的恒等映射延拓为 F 到 E 的一个映射）。
\item
  设 \(u\) 是从极小型空间 \(\mathbf{E}\) 到 Hausdorff 局部凸空间 \(\mathbf{F}\) 的一个连续线性映射。证明：\(u(\mathbf{E})\) 在 \(\mathbf{F}\) 中闭，且 \(u\) 是 \(\mathbf{E}\) 到 \(\mathbf{F}\) 的严格态射（利用 \(b\) 以及极小型空间的定义）。
\item
  设 F 是 Hausdorff 局部凸空间，M 是 F 的一个闭向量子空间。证明：如果存在 M 在 F 中的一个补 N，它是极小型的子空间，那么 N 是 M 在 F 中的拓扑补（利用 c)）。
\item
  设 M 是 Hausdorff 局部凸空间 F 中的一个极小型子空间；证明：对 F 的每个闭向量子空间 N，和 \(M + N\) 在 F 中闭（考察商空间 F/N 并利用 c)）。如果进而 N 也是极小型的，那么 \(M + N\) 是极小型的。
\end{enumerate}

\begin{enumerate}
\def\labelenumi{\arabic{enumi})}
\setcounter{enumi}{13}
\tightlist
\item
  设 \(\mathbf{E}\) 是 Hausdorff 局部凸空间，\(\mathbf{F}\) 是极小型的局部凸空间（习题 13）。
\end{enumerate}

\begin{enumerate}
\def\labelenumi{\alph{enumi})}
\item
  证明：如果 \(\mathbf{M}\) 是乘积空间 \(\mathrm{E} \times \mathrm{F}\) 的一个闭向量子空间，则它在 \(\mathrm{E}\) 上的投影在 \(\mathrm{E}\) 中闭（利用习题 13, e)）。
\item
  设 \(u\) 是 \(\mathbf{E}\) 到 \(\mathbf{F}\) 的一个线性映射。证明：如果 \(u\) 的图像在 \(\mathbf{E} \times \mathbf{F}\) 中闭，那么 \(u\) 连续（利用 \(a\) ）。
\item
  假设在 E 中每个闭向量子空间都有拓扑补（参见 V, 第 13 页）。证明：在 \(E \times F\) 中，每个闭向量子空间 M 都有拓扑补。（如果 \(N_{1}\) 是 M 在 E 上的投影，\(N_{2}\) 是 \(N_{1}\) 在 E 中的一个拓扑补，如果 \(P_{1} = M \cap F\) ，而 \(P_{2}\) 是 \(P_{1}\) 在 F 中的一个拓扑补，试利用 b) 证明 \(N_{2} + P_{2}\) 是 M 在 \(E \times F\) 中的一个拓扑补。）
\end{enumerate}

* 15) 设 E, F 是两个 Hausdorff 局部凸空间。我们称连续线性映射 \(u: \mathrm{E} \to \mathrm{F}\) 是线性真的，如果对每个 Hausdorff 局部凸空间 G 和 \(\mathrm{E} \times \mathrm{G}\) 的每个闭向量子空间 V，V 在映射 \(u \times 1_{\mathrm{G}}: \mathrm{E} \times \mathrm{G} \to \mathrm{F} \times \mathrm{G}\) 下的像是闭的。证明：这一条件等价于下述条件：\(u^{-1}(0)\) 是 E 的一个极小型子空间，且对 E 的每个闭向量子空间 W，集合 \(u(\mathbf{W})\) 在 F 中闭。（为证明第一个条件蕴含第二个，考察映射 \(v: \mathrm{E} \to \{0\}\) ，并给 E 装备拓扑 \(\sigma(\mathrm{E}, \mathrm{E}')\) ，使 E 浸入 \(\mathrm{E}'*\) 中且带拓扑 \(\sigma(\mathrm{E}'*, \mathrm{E}')\) ，然后取投影 \(v \times 1_{\mathrm{E}'*}: \mathrm{E} \times \mathrm{E}'* \to \mathrm{E}'*\) 下的像，这里所取的是 \(\mathrm{E} \times \mathrm{E}'*\) 中对角线 \(\Delta\) （\(\mathrm{E} \times \mathrm{E}\) 的对角线）的闭包。为证明第二个条件蕴含第一个，证明它蕴含：对拓扑 \(\sigma(\mathrm{E}, \mathrm{E}')\) 与 \(\sigma(\mathrm{F}, \mathrm{F}')\) 而言，映射 \(u\) 是严格态射，并利用习题 13, e)。) *

\begin{enumerate}
\def\labelenumi{\arabic{enumi})}
\setcounter{enumi}{15}
\tightlist
\item
  设 \(\mathbf{F}\) 是直线的乘积，\(\mathbf{C}\) 是 \(\mathbf{F}\) 中的一个闭凸集。
\end{enumerate}

\begin{enumerate}
\def\labelenumi{\alph{enumi})}
\item
  证明：存在 \(x_0 \in \mathbf{F}\) 、两个集合 I 与 J，以及一个拓扑同构 \(u\) （从 \(\mathbf{F}\) 到 \(\mathbf{R}^{\mathbf{I}} \times \mathbf{R}^{\mathbf{J}}\) 上的），使得 \(u(x_0 + \mathbf{C})\) 形如 \(\mathbf{R}^{\mathbf{I}} \times \mathbf{A}\) ，其中 \(\mathbf{A}\) 是 \(\mathbf{R}_+^{\mathbf{J}}\) 中的一个闭凸集。（注意我们有 \(\mathbf{F} = \mathbf{G}^*\) ，这里 \(\mathbf{F}\) 带拓扑 \(\sigma(\mathbf{G}^*, \mathbf{G})\) ；考察极集 \(C^\circ\) （\(C\) 在 \(G\) 中的极集）、\(G\) 中由 \(C^\circ\) 生成的向量子空间，以及该子空间的一个补。）
\item
  如果 C 不含任何仿射直线，那么 \((x, y) \mapsto x + y\) （从 \(C \times C\) 到 \(F\) 的映射）是真映射。
\item
  假设 C 是以 0 为顶点的锥，且由 F 的一致结构在 C 上诱导的一致结构是可度量化的。那么，如果集合 I 与 J、点 \(x_0\) 和映射 \(u\) 满足 \(a\) ) 的条件，则 I 是可数的，且存在 J 的一个可数子集 H，使得典范投影 \(p: \mathbf{R}^{\mathrm{I}} \times \mathbf{R}^{\mathrm{J}} \to \mathbf{R}^{\mathrm{I}} \times \mathbf{R}^{\mathrm{H}}\) 到 \(u(x_0 + \mathbf{C})\) 上的限制是一致子空间 \(u(x_0 + \mathbf{C})\) （\(\mathbf{R}^{\mathrm{I}} \times \mathbf{R}^{\mathrm{J}}\) 中的一致子空间）到一致子空间 \(p(u(x_0 + \mathbf{C}))\) （\(\mathbf{R}^{\mathrm{I}} \times \mathbf{R}^{\mathrm{H}}\) 中的一致子空间）上的同构。
\end{enumerate}

\begin{enumerate}
\def\labelenumi{\arabic{enumi})}
\setcounter{enumi}{16}
\tightlist
\item
  设 \(\mathbf{E}\) 是一个无限维向量空间。
\end{enumerate}

\begin{enumerate}
\def\labelenumi{\alph{enumi})}
\tightlist
\item
  证明：在 \(\mathrm{E}^*\) 中存在对拓扑 \(\sigma (\mathrm{E}^{*},\mathrm{E})\) 处处稠密的超平面。b) 如果 \(\mathbf{H}'\) 是这样一个超平面，证明：在 E 中，\(\neq \mathrm{E}\) 的、对拓扑 \(\sigma (\mathrm{E},\mathrm{H}')\) 处处稠密的线性子流形只有超平面。
\end{enumerate}

  ¶ 18) a) 在赋范空间 E 中，设 A 是一个闭凸集且 \(\neq\) E；证明函数 \(x\mapsto d(x,\complement\mathbf{A})\) 在 A 中是凹的（利用 A 是诸闭半空间之交这一事实）。b) 按下述方式归纳地定义闭凸集的一个序列 \(\mathbf{A}_{n}\subset\mathbf{R}^{n}\) ：\(\mathbf{A}_{1}=\mathbf{R}_{+}\) ；如果把 \(\mathbf{R}^{n+1}\) 等同于 \(\mathbf{R}^{n}\times\mathbf{R}\) ，那么 \(\mathbf{A}_{n+1}\) 是这样的偶 \((x,\xi)\) 所成之集：\(x\in\overset{\circ}{\mathbf{A}}_{n}\) 以及

\[
\xi \geqslant (d (x, \complement A _ {n})) ^ {- 1} + \| x \| ^ {2},
\]

其中 \(\|x\|\) 是欧氏范数。证明：\(A_{n+1}\) 没有任何形如 \(H \times R\) 的支撑超平面（这里 H 是 \(R^{n}\) 的一个超平面），且它的渐近锥是 \(\{0\} \times R_{+}\) 。c) 如果 \(p_{nm}\) 是典范投影 \(R^{m} \to R^{n}\) （把 \(R^{m}\) 等同于 \(R^{n} \times R^{m-n}\) ，对 \(m \geqslant n\) ），证明：当 \(R^{N}\) 等同于投影体系 ( \(R^{n}, p_{nm}\) ) 的投影极限时，诸 \(A_{n}\) 构成集合的一个投影体系，且 \(A = \varprojlim A_{n}\) 是 \(R^{N}\) 中一个非相对紧的闭凸集，它没有闭的支撑超平面，并且 \(C_{A} = \{0\}\) 。

\begin{enumerate}
\def\labelenumi{\arabic{enumi})}
\setcounter{enumi}{18}
\tightlist
\item
  \begin{enumerate}
  \def\labelenumii{\alph{enumii})}
  \tightlist
  \item
    设 A 是 E（直线的乘积）中一个非紧的闭凸集，且 \(C_{A} = \{0\}\) （习题 18）。证明：如果 \(B = \overline{A - A}\) ，且 M 是 \(A \cup (-A)\) 的闭凸包，那么 B 与 M 都含直线（利用 II, 第 85 页, 习题 16, b)）。
  \end{enumerate}
\end{enumerate}

\begin{enumerate}
\def\labelenumi{\alph{enumi})}
\setcounter{enumi}{1}
\tightlist
\item
  设 \(\mathrm{A}_1, \mathrm{A}_2\) 是 E 中两个闭凸集，使得 \(\mathrm{A}_1 + \mathrm{A}_2\) 闭，且 \(\mathrm{A}_1\) 、\(\mathrm{A}_2, \mathrm{A}_1 + \mathrm{A}_2\) 中没有一个含仿射直线。证明 \(\mathrm{C}_{\mathrm{A}_1 + \mathrm{A}_2} = \mathrm{C}_{\mathrm{A}_1} + \mathrm{C}_{\mathrm{A}_2}\) （利用 II, 第 85 页, 习题 16, b)）。c) 设 A 是 E 中不含任何仿射直线的一个闭凸集，而 \(\mathbf{M}_1, \dots, \mathbf{M}_n\) 是含于 A 中的闭凸集。如果 B 是 \(\bigcup_{i} \mathbf{M}_i\) 的凸包，证明 \(\overline{\mathbf{B}} = \mathbf{B} + \sum_{i} \mathbf{C}_{\mathbf{M}_i}\)
\end{enumerate}

以及 \(C_B = \sum_i C_{M_i}\) （同样的方法）。

¶ 20) 设 \(\mathrm{F} = \mathbf{R}^{(\mathrm{A})}, \mathrm{G} = \mathbf{R}^{\mathrm{A}}\) ，其中 \(\mathrm{A}\) 是任意无限集；假设 \(\mathrm{F}\) 与 \(\mathrm{G}\) 通过双线性型 \(\langle x, y \rangle = \sum_{\alpha \in \mathbf{A}} x(\alpha) y(\alpha)\) 处于分离对偶中。

\begin{enumerate}
\def\labelenumi{\alph{enumi})}
\item
  设 N 是 G 的一个加法子群；我们用 \(N^{*}\) 表示这样的 \(x \in F\) 所成的子群：\(\langle x, y \rangle\) 对所有 \(y \in N\) 都是整数；并用 \(N^{**}\) 表示这样的 \(z \in G\) 所成的子群：\(\langle x, z \rangle\) 对所有 \(x \in N^{*}\) 都是整数。如果 \(\overline{N}\) 是 N 对拓扑 \(\sigma(G, F)\) 的闭包，证明：\(N^{*}\) 在 F 中对 \(\sigma(F, G)\) 闭，且 \(N^{**} = \overline{N}\) （为确立最后一点，利用 GT, VII, § 1.3, 命题 6，把 N 投影到 G 的有限维坐标流形上）。
\item
  假设 \(\mathbf{A} = \mathbf{N}\) 。设 \(\mathbf{M}\) 是 \(\mathbf{F}\) 的一个对 \(\sigma(\mathbf{F},\mathbf{G})\) 闭的子群；证明：如果 \(\mathbf{V}\) 是含于 \(\mathbf{M}\) 中的最大向量子空间，那么 \(\mathbf{M}\) 是 \(\mathbf{V}\) 与一个闭子群 \(\mathbf{P}\) （它是一个具有可数基的自由 \(\mathbf{Z}\) -模）的拓扑直和。（把 \(\mathbf{F}\) 看作有限维向量子空间的递增序列 \((\mathrm{F}_n)\) 之并，并应用 GT, VII, § 1.2, 定理 2 与 § 1, 习题 7。）\(\mathbf{P}\) 是离散的（对由 \(\sigma(\mathbf{F},\mathbf{G})\) 诱导的拓扑而言），当且仅当 \(\mathbf{P}\) 是有限秩的。
\item
  由 a) 与 b) 推出：当 A = N 时，G 的每个闭子群（当 G 带乘积拓扑 \(\sigma(\mathbf{G}, \mathbf{F})\) 时）都可被拓扑群 G 的一个自同构变换为乘积 \(R^{I} \times Z^{J}\) ，其中 I 与 J 是 N 的两个无公共元素的集合。d) 在空间 \(E = R^{N}\) （带拓扑 \(\sigma(E, E^{*})\) ）中，证明子群 \(Z^{N}\) 是闭的且不含任何直线，尽管它不是自由 Z-模（A, VII, 第 59 页, 习题 8）；因此当 A 不可数时 b) 的结果不能推广。
\end{enumerate}

